\section{Struktur Dokumen HTML5}
HTML adalah fondasi struktur halaman web dan menjelaskan elemen apa saja yang ada di sebuah halaman. Dokumen HTML5 memiliki struktur standar seperti \code{doctype}, \code{head}, dan \code{body}. Pemahaman struktur ini penting agar browser dapat menampilkan konten dengan benar \cite{mdnHtmlIntro}.

Bagian \code{head} berisi metadata seperti judul halaman dan informasi viewport. Bagian \code{body} berisi seluruh konten yang terlihat oleh pengguna. Dengan struktur yang rapi, Anda dapat mengembangkan halaman lebih mudah dan terorganisir.

\begin{webcode}[language=HTML, caption={Struktur Dasar HTML5}]
<!doctype html>
<html lang="id">
  <head>
    <meta charset="utf-8" />
    <meta name="viewport" content="width=device-width, initial-scale=1" />
    <title>Halaman Pertama</title>
  </head>
  <body>
    <h1>Halo Web!</h1>
    <p>Ini adalah paragraf pertama.</p>
  </body>
</html>
\end{webcode}

